\phantomsection
\color{uniblue}\section*{\centering{\large{ПЕРЕЧЕНЬ СОКРАЩЕНИЙ}\color{white!0}<.>}}
\addcontentsline{toc}{section}{Перечень сокращений}
\color{black}
\begin{longtable}{>{\raggedright\arraybackslash}m{2cm}>{\raggedright\arraybackslash}m{0.5cm}>{\raggedright\arraybackslash}m{20cm}}
\endfirsthead\endhead\endfoot\endlastfoot
GOOSE & -- & Generic Object-Oriented Substation Event; \\
IRF & -- & Internal Relay Fault; \\
PPS & -- & Pulse Per Second; \\
АВ & -- & автоматический выключатель; \\
АОСН & -- & автоматика ограничения снижения напряжения; \\
АОСЧ & -- & автоматическое ограничение снижения частоты; \\
АПВН & -- & автоматическое повторное включение по напряжению; \\
АСУ~ТП & -- & автоматизированная система управления технологическими процессами; \\
АЦП & -- & аналого-цифровой преобразователь; \\
АЧР & -- & автоматическая частотная разгрузка; \\
БИ & -- & блок испытательный; \\
БНН & -- & блокировка при неисправности цепей напряжения; \\
БССЧ & -- & блокировка по скорости снижения частоты; \\
ВЭ & -- & выдвижной элемент (выкатная тележка, выкатной элемент); \\
ЗМН & -- & защита минимального напряжения; \\
ЗН & -- & заземляющий нож (разъединитель); \\
ЗФР & -- & защита от феррорезонанса; \\
ИО & -- & измерительный орган; \\
ИЧМ & -- & интерфейс человек-машина; \\
КА & -- & коммутационный аппарат; \\
КНН & -- & контроль наличия напряжения; \\
КОН & -- & контроль отсутствия напряжения; \\
КП & -- & контроллер присоединения; \\
КПОН & -- & комбинированный пусковой орган напряжения; \\
КС & -- & контроль синхронизма; \\
КССШ & -- & контроль смежной секции шин; \\
ЛФ & -- & логика формирования; \\
МУВ & -- & матрица управляющих воздействий; \\
НКУ & -- & низковольтное комплектное устройство; \\
НП & -- & нулевая последовательность; \\
ОБР & -- & оперативная блокировка разъединителей; \\
ОВ & -- & оперативный вывод; \\
ОЗУ & -- & оперативное запоминающее устройство; \\
ОН & -- & орган направления / отключение нагрузки; \\
ОП & -- & обратная последовательность; \\
ОТ & -- & оперативный ток; \\
ПА & -- & противоаварийная автоматика; \\
ПДС & -- & преобразователь дискретных сигналов; \\
ПЗУ & -- & постоянное запоминающее устройство; \\
ПО & -- & пусковой орган / программное обеспечение; \\
ПП & -- & прямая последовательность; \\
ПС & -- & предупредительная сигнализация; \\
РАС & -- & регистратор аварийных событий; \\
РЗ & -- & релейная защита; \\
РЗА & -- & релейная защита и автоматика; \\
РЭ & -- & руководство по эксплуатации; \\
САЧР & -- & специальная очередь автоматической частотной разгрузки; \\
СВ & -- & секционный выключатель; \\
СЗЗ & -- & сигнализация замыкания на землю; \\
СКРМ & -- & средства компенсации реактивной мощности; \\
СС & -- & сборка сигналов; \\
ТН & -- & трансформатор напряжения; \\
ТТ & -- & трансформатор тока; \\
УВ & -- & управление выключателем / управляющее воздействие; \\
УВЭ & -- & управление выкатным элементом; \\
УЗН & -- & управление заземляющим ножом (разъединителем); \\
УОН & -- & устройство отключения нагрузки; \\
ФК & -- & функциональная клавиша; \\
ФСУ & -- & функциональная схема устройства; \\
ЦВН & -- & централизованное включение нагрузки; \\
ЦН & -- & цепи напряжения; \\
ЦП & -- & центральный процессор; \\
ЧАПВ & -- & частотное автоматическое повторное включение. \\
\end{longtable}